\subsection{4. Dinamika Pelaksanaan UUD NRI Tahun 1945}
Pelaksanaan UUD NRI Tahun 1945 telah mengalami dinamika dan beberapa kali pergantian akibat perkembangan politik nasional:

\subsubsection{A. Periode UUD NRI 1945 (18 Agustus 1945 -- 27 Desember 1949)}
\begin{itemize}
    \item Pemerintahan awal fokus menghadapi ancaman kembalinya Sekutu dan Belanda.
    \item Terjadi perubahan sistem pemerintahan dari \textbf{Presidensial} menjadi \textbf{Parlementer} melalui Maklumat Pemerintah tanggal 14 November 1945 (munculnya Kabinet Sjahrir, Hatta, dll.).
\end{itemize}

\subsubsection{B. Periode Konstitusi RIS (27 Desember 1949 -- 17 Agustus 1950)}
\begin{itemize}
    \item Berdasarkan hasil Konferensi Meja Bundar (KMB).
    \item \textbf{Bentuk Negara}: Serikat / Federasi.
    \item \textbf{Bentuk Pemerintahan}: Republik.
    \item \textbf{Sistem Pemerintahan}: Parlementer.
\end{itemize}

\subsubsection{C. Periode UUDS 1950 (17 Agustus 1950 -- 5 Juli 1959)}
\begin{itemize}
    \item Bentuk negara kembali menjadi \textbf{Kesatuan}.
    \item \textbf{Sistem Pemerintahan}: Parlementer.
    \item \textbf{Dampak Negatif Masa UUDS 1950}:
    \begin{itemize}
        \item Sering terjadi pergantian kabinet (7 kali berganti kabinet dalam 9 tahun), menyebabkan program jangka panjang terhambat.
        \item Timbul ketidakharmonisan dan pemberontakan di daerah (PRRI/Permesta).
        \item Kebuntuan Badan Konstituante hasil Pemilu 1955 dalam menyusun UUD baru.
    \end{itemize}
\end{itemize}

\subsubsection{D. Periode UUD 1945 Masa Orde Lama (5 Juli 1959 -- 1965/1966)}
Diawali dengan \textbf{Dekrit Presiden 5 Juli 1959} yang berisi:
\begin{enumerate}
    \item Pembubaran Badan Konstituante.
    \item Berlakunya kembali UUD 1945 dan tidak berlakunya UUDS 1950.
    \item Pembentukan MPRS dan DPAS.
\end{enumerate}

\paragraph{Bentuk Penyimpangan Masa Orde Lama (Demokrasi Terpimpin)}
\begin{itemize}
    \item Presiden mengeluarkan \textit{Penetapan Presiden (Penpres)} tanpa persetujuan DPR.
    \item Pengangkatan Presiden Sukarno sebagai Presiden Seumur Hidup (TAP MPRS No. III/MPRS/1963).
    \item Pembubaran DPR hasil Pemilu 1955 oleh Presiden dan diganti dengan DPR-Gotong Royong (DPR-GR).
    \item Kekuasaan terpusat sangat dominan pada Presiden (\textit{Executive Heavy}).
\end{itemize}

\subsubsection{E. Periode UUD 1945 Masa Orde Baru (1966 -- 1998)}
Tekad awal: Melaksanakan Pancasila dan UUD 1945 secara murni dan konsekuen.

\paragraph{Bentuk Penyimpangan Masa Orde Baru}
\begin{itemize}
    \item Penyelenggaraan ekonomi kurang memperhatikan Pasal 33 UUD 1945 (terjadi monopoli dan pemusatan ekonomi).
    \item Kekuasaan Presiden sangat dominan (\textit{Executive Heavy}).
    \item Pengekangan kebebasan pers, berpendapat, dan berorganisasi.
    \item Maraknya praktik \textbf{KKN (Korupsi, Kolusi, dan Nepotisme)}.
    \item Penegakan hukum lemah dan rentan diintervensi penguasa.
\end{itemize}


\subsection{5. Pergantian dan Perubahan UUD NRI Tahun 1945}

\subsubsection{Sifat Konstitusi}
\begin{table}[h!]
\centering
\begin{tabular}{|l|p{10cm}|}
\hline
\textbf{Sifat Konstitusi} & \textbf{Penjelasan} \\ \hline
\textbf{Rigid (Kaku)} & Memerlukan tata cara dan persyaratan khusus yang sulit untuk diubah (contoh: UUD NRI 1945 Pasal 37). \\ \hline
\textbf{Flexible (Luwes)} & Cara perubahannya tidak berbeda dengan cara pembuatan undang-undang biasa. \\ \hline
\end{tabular}
\end{table}

\subsubsection{Perubahan UUD NRI 1945 Era Reformasi (1999 -- 2002)}
Amandemen dilakukan sebanyak \textbf{4 kali} dalam Sidang Umum/Tahunan MPR (1999, 2000, 2001, dan 2002) untuk membatasi kekuasaan \textit{executive heavy}, memperkuat perlindungan HAM, serta menegaskan otonomi daerah.

\paragraph{Lima Kesepakatan Dasar dalam Amandemen UUD NRI 1945}
\begin{enumerate}
    \item Tidak mengubah Pembukaan UUD NRI Tahun 1945.
    \item Tetap mempertahankan Negara Kesatuan Republik Indonesia (NKRI).
    \item Mempertegas sistem pemerintahan presidensial.
    \item Penjelasan UUD NRI Tahun 1945 yang memuat hal-hal normatif dimasukkan ke dalam pasal-pasal.
    \item Perubahan UUD dilakukan dengan cara \textit{addendum} (mempertahankan naskah asli dan melampirkan naskah perubahan di bagian akhir).
\end{enumerate}

\paragraph{Wawasan Kewarganegaraan: Pengertian Amendemen}
\textit{Amendemen} berasal dari bahasa Inggris (\textit{to amend} = menyempurnakan/memperbaiki). Amendemen UUD NRI 1945 \textbf{bukanlah penggantian total} konstitusi, melainkan melengkapi, menyempurnakan, atau membatasi klausul pasal-pasal tertentu tanpa mengubah naskah dasarnya, agar sesuai dengan tuntutan zaman dan keadilan masyarakat.

\subsubsection{Perubahan UUD NRI Tahun 1945 (Sebelum dan Setelah Amendemen)}

\begin{table}[h!]
\centering
\small
\begin{tabular}{|p{3cm}|p{6cm}|p{6cm}|}
\hline
\textbf{Perihal} & \textbf{Sebelum Amendemen} & \textbf{Setelah Amendemen} \\ \hline
\textbf{Sistematika Batang Tubuh / Pasal} & 16 bab, 37 pasal, 49 ayat, 4 pasal Aturan Peralihan, dan 2 ayat Aturan Tambahan. & 21 bab, 73 pasal (keseluruhan) dan 170 ayat, 3 pasal Aturan Peralihan, dan 2 ayat Aturan Tambahan. \\ \hline
\textbf{Kedaulatan Rakyat} & Kedaulatan adalah di tangan rakyat dan dilakukan sepenuhnya oleh Majelis Permusyawaratan Rakyat. & Kedaulatan berada di tangan rakyat dan dilaksanakan menurut UUD. \\ \hline
\textbf{Kekuasaan Presiden (Eksekutif)} & Presiden memiliki kekuasaan yang besar (\textit{concentration of power and responsibility upon the president}). & Prinsip saling mengawasi dan mengimbangi (\textit{checks and balances}). \\ \cline{2-3}
 & Presiden dipilih oleh MPR. & Presiden dipilih langsung oleh rakyat melalui pemilihan Presiden dan Wakil Presiden. \\ \hline
\textbf{Kedudukan Lembaga Negara} & MPR sebagai lembaga tertinggi negara. & Setiap lembaga negara sejajar kedudukannya di bawah UUD NRI Tahun 1945. \\ \hline
\textbf{Lembaga Negara} & Keberadaan Dewan Pertimbangan Agung. & Pembubaran DPA dan pembentukan lembaga baru seperti Mahkamah Konstitusi, Komisi Yudisial, dan Dewan Perwakilan Daerah (DPD). \\ \cline{2-3}
 & Keanggotaan MPR: berasal dari anggota DPR ditambah Utusan Daerah dan Utusan Golongan. & Keanggotaan MPR: berasal dari anggota DPR ditambah Dewan Perwakilan Daerah (DPD). \\ \cline{2-3}
 & Anggota DPR dipilih melalui pemilu ditambah anggota TNI/Polri yang diangkat. & Anggota DPR dipilih melalui pemilu. \\ \cline{2-3}
 & Bank sentral tidak disebut secara eksplisit dalam UUD. & Bank sentral disebut secara eksplisit dalam UUD. \\ \hline
\textbf{Pejabat Negara} & Belum mengatur mekanisme pengangkatan dan pemberhentian para pejabat negara. & Mengatur mekanisme pengangkatan dan pemberhentian para pejabat negara. \\ \hline
\textbf{Prinsip Negara Hukum} & Indonesia ialah negara yang berdasar atas hukum (terdapat pada penjelasan UUD). & Mempertegas prinsip negara hukum dengan menempatkan kekuasaan kehakiman sebagai kekuasaan yang merdeka, penghormatan kepada HAM serta kekuasaan yang dijalankan atas prinsip \textit{due process of law}. \\ \hline
\textbf{Sistem Konstitusional} & Pemusatan kekuasaan pada Presiden atau eksekutif (\textit{executive heavy}). & Sistem konstitusional berdasarkan perimbangan kekuasaan (\textit{checks and balances}), yaitu setiap kekuasaan dibatasi oleh undang-undang berdasarkan fungsi masing-masing. \\ \hline
\textbf{Hak Asasi Manusia (HAM)} & HAM tidak dijabarkan secara rinci dan tegas pada pasal-pasal UUD. & Pasal mengenai HAM dijabarkan secara tegas (Pasal 28 sampai dengan 34 UUD). \\ \hline
\end{tabular}
\end{table}